\label{sec:results}
% <- T8; ablations from paper/results_digest.md; identifiability LEADS.
% CAVEAT (state once, applies to every subsection): all T8 numbers are N=1
% (single seed=42 run per cell); no confidence intervals are statistically
% defensible at N=1 (STATS_PROTOCOL Stage-1-lite scoping). T8 as executed is
% a single-airfoil (NACA 2412) ablation matrix plus two extra single-airfoil
% cells (0012, 4415), NOT the pre-registered ~20-section NACA panel
% (configs/experiments/panel_naca.yaml does not exist). Source:
% paper/results_digest.md preamble; experiments/results/t8/ANALYSIS.md \S0.

Every number in this section traces to a \texttt{result.json} or
\texttt{diagnostics.json} artifact under \pth{experiments/results/t8/}
(paper/results\_digest.md; per-cell citation in
\pth{experiments/results/t8/ANALYSIS.md}). All values are $N{=}1$
(single \texttt{seed=42} run per cell); no confidence intervals are reported
because none are statistically defensible at $N{=}1$, per
\pth{docs/STATS_PROTOCOL.md}'s Stage-1-lite scoping. The $n$-sweep,
station-selection, and other single-factor \emph{ablation} tables in this
section are a single-airfoil (NACA 2412) matrix plus two extra
single-airfoil cells (0012, 4415). These are deliberately narrow,
$N{=}1$-per-cell comparisons designed to isolate one factor at a time, not
a panel study.
Panel-level convergence-\emph{fraction} language is deliberately avoided in
those subsections for that reason. \S\ref{sec:h1-panel} is the one
exception: it reports the pre-registered $\approx$20-section NACA panel
(\pth{configs/experiments/panel_naca/*.yaml}), run subsequently at the
winning configuration identified by this section's ablations, and its
Wilson-interval convergence-fraction claim is reported there on its own
terms.

\paragraph{Station-addressing provenance.} The ablation matrix in this
section was run with target stations addressed by panel-node index, which was
the implementation at the time. Addressing was subsequently changed to
$(\text{surface}, x/c)$ with interpolation between the two bracketing nodes,
because a node index stops denoting the same physical location once the
geometry moves, which is precisely what an inverse solve does
(\S\ref{sec:newton-system}). The change moves the T7 figures quoted in
\S\ref{sec:falsifiable} from $1.08\times10^{-11}$ in 4 iterations to
$2.75\times10^{-11}$ in 6, both far inside the $10^{-4}$/9-iteration gate,
and leaves the QR-selected submap conditioning unchanged at 90.3. Every cell
in this section shares the earlier addressing, so the within-matrix
comparisons that the ablations rest on are unaffected; what must not be done
is to read a single number from this section alongside a single number from
\S\ref{sec:falsifiable} as though they came from one run. Where that
distinction matters, it is stated at the point of use.

\subsection{Station-selection identifiability: the central methodological result}
\label{sec:ablations}

This finding leads the ablation discussion, ahead of the flow-solve count,
because it is the most novel, most precisely quantified, and most
consequential result in T8 (benchmarks/REPORT.md, ``Recommendations for
P1''). The square extended Newton system (\S\ref{sec:newton-system}) demands
exactly $M$ target-station equations. With evenly-spaced stations, between-
station Cp information is discarded, and the resulting $16\times16$ station-
sensitivity submap can be badly ill-conditioned. Table~\ref{tab:identifiability}
reports the two configurations that isolate this mechanism.

\begin{table}[h]
\centering
\caption{Station-selection ablation, the paper's central methodological
result. Both cells report \texttt{converged: true} and drive the residual to
the solver floor, yet their recovery errors differ by eight orders of
magnitude: only the QR-pivoted selection recovers the generating design, and
the evenly-spaced selection converges to a spurious root while its station
submap conditions at $3\times10^{12}$. Source:
\pth{experiments/results/t8/ANALYSIS.md}~\S4.}
\label{tab:identifiability}
\begin{tabular}{lccc}
\toprule
Selection & $\mathrm{err\_all\_inf}$ & Submap $\mathrm{cond}$ & Newton \texttt{converged} flag \\
\midrule
\texttt{qr\_pivot} (winning config) & $\mathbf{1.08\times10^{-11}}$ & 90.3 & \texttt{true} \\
\texttt{even} (evenly-spaced) & $\mathbf{1.43\times10^{-3}}$ (FAILS $10^{-4}$ gate) & $3.47\times10^{12}$ & \texttt{true} \\
\bottomrule
\end{tabular}
\end{table}

Both cells' Newton loops report their own internal \texttt{converged: true};
the residual reaches the solver's floor in both cases. \textbf{Only
\texttt{qr\_pivot} recovers the correct design.} ``Converged'' and
``correct'' are different claims, invisible unless $\mathrm{err\_all\_inf}$
is checked against the true coefficients rather than the residual norm
alone (experiments/results/t8/ANALYSIS.md~\S0, Deviation 5). The mechanism:
with evenly-spaced stations, near-null coefficient directions exist whose
motion changes the sampled Cp by less than floating-point precision.
Newton converges to \emph{any} of several exact roots of the (badly-
conditioned) discrete system, a structural non-uniqueness of the sampled
problem, not a numerical accident. This reproduces and sharpens a pre-T8
finding (\pth{docs/triz/T7-identifiability.md}: cond $\sim10^7$, error
$1.4\times10^{-3}$--$6\times10^{-2}$ with evenly-spaced stations) with a
controlled ablation cell. Two contributing null-direction families were
identified and eliminated \emph{before} the station-selection fix, each
verified empirically (\pth{docs/triz/T7-identifiability.md}): a
camber--$\alpha$ equivalence (freeing $\alpha$ lets a slightly different
camber at a slightly different $\alpha$ interpolate the same samples, so
$\alpha$ is fixed for self-consistency tests), and a prescribed-LE mismatch
(excluding LE stations while leaving the LE-dominant coefficients free
leaves $A_0$ unobservable, which is why \texttt{le\_treatment=prescribed}
\emph{fixes} $A_{u,0}$/$A_{l,0}$ rather than merely dropping LE-region
target rows).

The resolution keeps the system square, with no least-squares relaxation
and no reintroduced optimisation, and instead makes each of the $M$
equations maximally informative. The T4 pre-solve (\S\ref{sec:presolve})
already builds the full $(200\times n_A)$ sensitivity matrix for
initialisation; reusing it to \emph{select} the $M$ target stations by QR
column pivoting chooses the rows that maximise the resulting square
submap's conditioning. The measured effect: submap condition
$10^7\to90.3$; recovery error $1.4\times10^{-3}\to1.08\times10^{-11}$.
Full-curve sensitivity information enters only through the
(already-required) linear pre-solve, and the nonlinear Newton solve stays
square and determined throughout (\pth{docs/triz/T7-identifiability.md}).
Sensitivity-optimal station selection is therefore the identifiability
guard for determined inverse systems, a methodological contribution of this
paper beyond the original dossier hypothesis: a practical answer to the
non-uniqueness caveat inverse problems inherit from Lighthill's and
Volpe \& Melnik's well-posedness constraints
(\S\ref{sec:presolve}; \dossierref{\S7.10}),
without abandoning the root-finding architecture for a least-squares
relaxation.

Figure~\ref{fig:station-selection} summarises the ablation of
Table~\ref{tab:identifiability}: the same solver, the same target and the
same convergence flag produce recovery errors eight orders of magnitude
apart, separated only by which stations were sampled.

\begin{figure}[h]
\centering
\includegraphics[width=0.75\textwidth]{station_selection.png}
\caption{Station-selection identifiability at a glance: recovery error, QR-pivoted vs.\ evenly-spaced target stations, annotated with submap conditioning. Source: \pth{experiments/results/t8/figures/station_selection.png}.}
\label{fig:station-selection}
\end{figure}

\subsection{Where the stations are placed: the prescribed leading edge}
\label{sec:le-identifiability}

\S\ref{sec:ablations} records one half of a requirement: excluding
leading-edge target stations while leaving the LE-dominant coefficients free
leaves $A_0$ unobservable, which is why \texttt{le\_treatment=prescribed}
fixes $A_{u,0}$ and $A_{l,0}$ rather than merely dropping LE-region rows. The
other half is the converse, and it was found while extending the method to
arbitrary user-supplied targets: once the leading edge \emph{is} prescribed,
the LE-region target stations must in turn be excluded. The two conditions are
not independent choices. Prescribing the nose without restricting the
candidate stations, or restricting the stations without prescribing the nose,
each leaves a near-null direction in the sampled system; only the pair is
identifiable.

Under the prescribed-LE treatment (FM-3, \S\ref{sec:fm}), $A_{u,0}$ and
$A_{l,0}$ are given rather than solved, so a target row placed inside the
prescribed region constrains pressure over a piece of surface that cannot
move. The effect on sampling is not a uniform loss of information, and
Figure~\ref{fig:le-identifiability}(a) is worth reading carefully on this
point. Across the prescribed region the norm of
$\partial C_p/\partial A_{\text{free}}$ spans nearly two orders of magnitude,
from $0.054$ at the nose itself to $3.90$ just aft of it, against a median of
$1.47$ over the rest of the chord. The region therefore contains both the
single least informative station on the surface and some of the most
sensitive. That combination is what makes unrestricted QR selection a poor
choice here rather than merely a neutral one: the pivoting is drawn to the
high-sensitivity nose stations, and the rows it selects there are strongly
mutually correlated, because a small neighbourhood of a fixed nose produces
nearly the same pressure response for every free coefficient. The extended
system stays formally square and full-rank while becoming numerically
near-dependent, which is the same underlying condition as
Table~\ref{tab:identifiability} reached from the opposite direction.

Restricting the QR candidate set to $x \geq$ the prescribed-LE fraction, as
Algorithm~\ref{alg:qrpivot} specifies, removes the degeneracy.
Table~\ref{tab:le-identifiability} reports the effect on a self-consistency
target, holding everything else fixed.

\begin{table}[h]
\centering
\caption{Target-station placement under a prescribed leading edge. Both cells
start from the same perturbed geometry and differ only in whether the QR
candidate set is restricted to $x \geq$ the prescribed-LE fraction. With the
leading-edge region left in the candidate set the solve does not converge at
all within the 50-iteration budget, despite driving its residual to
$1.9\times10^{-10}$. Surface offset is quoted in millichords, against the 1
millichord that \S\ref{sec:formulation}'s fit gate allows a CST
representation itself. Source:
\pth{experiments/results/le_stations/result.json}.}
\label{tab:le-identifiability}
\footnotesize
\begin{tabularx}{\textwidth}{@{}L c c@{}}
\toprule
Quantity & LE stations in candidate set & LE stations excluded \\
\midrule
Candidate stations          & 200                   & 169 \\
Submap condition            & $31.7$                & $19.9$ \\
Converged                   & \textbf{no}           & \textbf{yes} \\
Newton iterations           & 50 (budget exhausted) & \textbf{7} \\
Final residual              & $1.89\times10^{-10}$  & $8.21\times10^{-11}$ \\
$\norm{A-A^*}_\infty$ (free)& $6.97\times10^{-1}$   & $\mathbf{4.10\times10^{-3}}$ \\
Max surface offset (mc)     & $11.74$               & $\mathbf{0.127}$ \\
\bottomrule
\end{tabularx}
\end{table}

The residual is not what separates the two columns. The unrestricted cell
drives its residual to $1.9\times10^{-10}$, within an order of magnitude of
the restricted cell's, and still fails to converge inside a 50-iteration
budget while ending $11.7$ millichords from the target surface. Restricting
the candidate set converges in 7 iterations to $0.127$ millichords, a
ninety-fold improvement in the recovered shape, with the submap conditioning
improving only from $31.7$ to $19.9$. A conditioning number that moves by
less than a factor of two therefore does not bound the effect: what matters
is whether the equations are spent on a part of the surface the solver is
allowed to move.

This reinforces rather than repeats \S\ref{sec:ablations}. There, the
degeneracy came from discarding between-station information across the whole
chord, and both configurations converged, one of them to a wrong root. Here
it comes from spending equations on a region that has been removed from the
design space, and the failure is visible as non-convergence rather than as a
confident wrong answer. Neither is detectable from the residual norm alone.

\begin{figure}[htbp]
\centering
\includegraphics[width=\textwidth]{fig_le_identifiability.png}
\caption{Why target stations must be kept out of a prescribed leading edge.
(a)~the sensitivity of surface pressure to the free coefficients across the
chord, with the prescribed region shaded: it does not vanish there but varies
by nearly two orders of magnitude, holding both the least and some of the most
sensitive stations; (b)~the stations chosen by unrestricted
QR against those chosen from the restricted candidate set, with the resulting
submap conditioning; (c)~the effect on the solve, plotted relative to the
unrestricted case on a logarithmic axis. Values as in
Table~\ref{tab:le-identifiability}.}
\label{fig:le-identifiability}
\end{figure}

\subsection{Bernstein order $n$: the FM-2 conditioning cliff}

Table~\ref{tab:n-sweep} sweeps Bernstein order $n$ (per side) from 4 to 12,
reporting the three ADR-0004-distinct conditioning quantities alongside
iteration count and recovery error. These are distinct matrices, plotted
and reported separately per \S\ref{sec:presolve}'s reporting rule.

\begin{table}[h]
\centering
\caption{Bernstein order sweep, $n=4$ to $12$ per side. Iteration count is
flat at 3 to 4 through $n=10$ and then jumps to 21 at $n=12$, a cliff rather
than a gradual degradation, while all three conditioning measures climb
monotonically throughout. Source:
\pth{experiments/results/t8/ANALYSIS.md}~\S4.}
\label{tab:n-sweep}
\begin{tabular}{lccccc}
\toprule
$n$ & Iterations & $\mathrm{err\_all\_inf}$ & Submap $\mathrm{cond}$ & $\mathrm{cond}(J)$, D-2, it.\,0 & T2 Gram/fit $\mathrm{cond}$ \\
\midrule
4  & 3 & $8.08\times10^{-12}$ & 10.6 & $6.46\times10^7$ & 498 \\
6  & 4 & $1.59\times10^{-10}$ & 20.4 & $8.46\times10^7$ & 7{,}327 \\
8  & 4 & $1.08\times10^{-11}$ & 90.3 & $1.57\times10^8$ & 108{,}538 \\
10 & 4 & $1.46\times10^{-10}$ & 592  & $4.52\times10^8$ & 1{,}626{,}188 \\
\textbf{12} & \textbf{21} (FAILS 9-iter gate) & $6.31\times10^{-10}$ & 1{,}144 & $2.07\times10^9$ & 24{,}579{,}174 \\
\bottomrule
\end{tabular}
\end{table}

Iteration count is flat at 3--4 through $n=10$, then jumps to 21 at $n=12$.
This is a genuine cliff, not gradual degradation, though it tracks (without
being identical to) all three independently-measured conditioning series,
which climb monotonically with $n$ throughout the sweep (roughly
$4^n$-like growth in the T2 Gram/fit condition, consistent with the a
priori FM-2 hypothesis, \dossierref{\S4}
FM-2). $n=12$ still eventually reaches the accuracy gate
($6.31\times10^{-10}<10^{-4}$); it fails only the iteration-count budget,
which is itself informative: FM-2 degrades convergence \emph{speed} before
it degrades final accuracy. Figure~\ref{fig:nsweep}.

\begin{figure}[h]
\centering
\includegraphics[width=0.85\textwidth]{n_sweep.png}
\caption{Iterations and $\mathrm{err\_all\_inf}$ vs.\ Bernstein order $n$, with the three ADR-0004-distinct conditioning series plotted separately. Source: \pth{experiments/results/t8/figures/n_sweep.png}.}
\label{fig:nsweep}
\end{figure}

\subsection{FM-1, FM-4, and other single-factor ablations}

Table~\ref{tab:other-ablations} reports the remaining single-factor
ablations against the \texttt{t8\_n08\_baseline} cell.

\begin{table}[h]
\centering
\caption{Single-factor ablations vs.\ baseline. Source: paper/results\_digest.md~\S5.2.}
\label{tab:other-ablations}
\footnotesize
\begin{tabularx}{\textwidth}{@{}l c c c L@{}}
\toprule
Factor & Cell & Iterations & $\mathrm{err\_all\_inf}$ & Notes \\
\midrule
Transition free vs.\ forced & \texttt{t8\_transition\_free} & 4 & $1.28\times10^{-11}$ & no material difference (FM-4 mitigation cost is negligible once applied) \\
$\alpha$ free vs.\ fixed & \texttt{t8\_alpha\_free} & 5 & $2.31\times10^{-9}$ & 1 extra iteration; consistent with the camber--$\alpha$ null direction (\S\ref{sec:ablations}) \\
Init: perturbed vs.\ presolve & \texttt{t8\_init\_perturbed} & 5 & $2.12\times10^{-11}$ & 43 vs.\ 120 flow-solve-equiv (\S\ref{sec:h2}); presolve buys efficiency, not correctness \\
LE treatment: none vs.\ prescribed & \texttt{t8\_le\_none} & 4 & $6.50\times10^{-10}$ & submap cond roughly doubles (195.6 vs.\ 90.3), still passes gate by 6 orders of magnitude \\
DOF over-determined ($+1$) & \texttt{t8\_dof\_over} & 0 (clean fail) & n/a & Eq.~\eqref{eq:dof} violated; \texttt{"M+K=17 but n\_A\_free+n\_alpha=16"} \\
DOF under-determined ($-1$) & \texttt{t8\_dof\_under} & 0 (clean fail) & n/a & Eq.~\eqref{eq:dof} violated; \texttt{"M+K=15 but n\_A\_free+n\_alpha=16"} \\
\bottomrule
\end{tabularx}
\end{table}

The two DOF-mis-squared cells fail at construction time with an explicit,
correct error string and zero Newton iterations spent. The FM-1 guard
(Eq.~\eqref{eq:dof}) catches deliberate over/under-determination cleanly:
no silent corruption, no wrong answer (experiments/results/t8/ANALYSIS.md~\S4).

\subsection{H1: NACA panel generalization}
\label{sec:h1-panel}

Table~\ref{tab:airfoils} adds two additional single-airfoil cells
(0012, 4415) to the 2412 baseline reported throughout
\S\ref{sec:falsifiable}; that 3-airfoil table alone is $N{=}1$ per section
and explicitly not a generalization claim
(paper/results\_digest.md~\S5.2). It was subsequently superseded by the
pre-registered H1 evaluation itself, run post-review after the T8 ablation
matrix closed (experiments/results/t8/ANALYSIS.md, ``H1 addendum''):
the winning configuration (\S\ref{sec:falsifiable}) applied to a 20-section
NACA panel spanning 9--25\% thickness and 0--6\% camber across both 4- and
5-digit families (\pth{configs/experiments/panel_naca/*.yaml},
\texttt{seed=42}).

\begin{table}[h]
\centering
\caption{Three-airfoil detail (baseline + 2 extras). Source: paper/results\_digest.md~\S5.2.}
\label{tab:airfoils}
\begin{tabular}{lcccc}
\toprule
Airfoil & Iterations & $\mathrm{err\_all\_inf}$ & Submap $\mathrm{cond}$ & Model gap \\
\midrule
2412 (baseline) & 4 & $1.08\times10^{-11}$ & 90.3 & 0.110 \\
0012 & 4 & $3.15\times10^{-11}$ & 83.9 & 0.083 \\
4415 & 3 & $4.35\times10^{-11}$ & 117.9 & 0.163 \\
\bottomrule
\end{tabular}
\end{table}

\paragraph{Result: 18/18 generable sections recovered, at $\mathrm{err\_free\_inf}\leq1.51\times10^{-10}$
in 3--7 Newton iterations each.} Two of the 20 pre-registered sections
produced no result and are excluded under the pre-registered \S6-class
exclusion rules (experiments/results/t8/ANALYSIS.md, ``H1 addendum,''
Deviation~11), not counted as failures:
For \texttt{panel\_0006}, the forced-trip \emph{target} generation itself
fails to converge for this 6\%-thin section at a 5\% trip location, so no
inverse problem was ever posed (a direct-solve non-convergence exclusion,
exactly the pre-registered rule's intended case, not an architecture
failure). For \texttt{panel\_44012}, the vendor NACA 5-digit geometry
generator does not implement the 44XXX mean-line family, a pre-registration
defect (the target list included a section outside the implemented geometry
generator), not a solver failure. Figure~\ref{fig:h1-gallery} shows 6
representative recovered sections spanning this range: thin symmetric
(0009), thick symmetric (0025), the cambered baseline (2412), a thick
cambered section (4415), a standard (non-reflexed) 5-digit-camber NACA5 (23012), and a high-camber
section (6412). Target and recovered outlines are visually
indistinguishable at this scale, consistent with the $\leq1.51\times10^{-10}$
coefficient-level agreement.

\begin{figure}[h]
\centering
\includegraphics[width=0.95\textwidth]{paper/fig_h1_panel_gallery.png}
\caption{H1 panel gallery: 6 representative recovered sections (thin/thick/cambered/NACA5), target vs.\ recovered geometry, annotated with each cell's $\mathrm{err\_all\_inf}$ and iteration count from its \texttt{result.json}. Source: \pth{experiments/results/t8/figures/paper/fig_h1_panel_gallery.png} (regenerated by \pth{src/cins/benchmarks/paper_figures.py::fig_h1_panel_gallery}); scalar annotations from \pth{experiments/results/t8/panel_*/result.json}.}
\label{fig:h1-gallery}
\end{figure}

\paragraph{Interval-arithmetic limits of the panel size (the pre-registered
criterion still fails, on a technicality).} At $n_{\text{eff}}=18$
successes out of 18 eligible trials, the Wilson 95\% lower bound is
$\mathbf{0.824}$, \emph{below} the pre-registered criterion (Wilson LB
$\geq0.9$), even though every eligible section succeeded. This is not a
solver shortfall. A Wilson lower bound of 0.9 is mathematically
unattainable at $n=18$ or $n=20$ regardless of outcome (the smallest
all-success sample whose Wilson lower bound clears 0.9 is $n=35$, at
$0.901$; $n=29$ reaches only $0.883$): the pre-registered target and the
pre-registered panel size are mutually inconsistent, a pre-registration
defect (experiments/results/t8/ANALYSIS.md, Deviation~12). Treating
the excluded \texttt{panel\_0006} as a failure rather than an exclusion
instead gives 18/19 (0.947), Wilson LB 0.754, still short of 0.9, for the
same $n$-too-small reason. On the NACA panel alone, the architecture
converges and recovers the true coefficients on every panel section it
could be legitimately posed on, and the formal pre-registered statistical
criterion fails on interval arithmetic, not on solver performance.
\S\ref{sec:h1-uiuc-panel} reports the pre-registered UIUC extension, sized
to resolve the interval-arithmetic limit, and gives the combined H1
verdict across both panels.

\subsection{H1: UIUC panel extension}
\label{sec:h1-uiuc-panel}

Table~\ref{tab:uiuc-outcome} and Table~\ref{tab:uiuc-criteria} report the
pre-registered UIUC panel (STATS\_PROTOCOL~\S3.3). Sections were drawn from
the 123-file UIUC corpus, stratified by thickness and camber decile, after
a pre-filter removed 6 sections whose $n{=}8$ CST fit exceeds the T2 gate
(\pth{configs/experiments/panel_uiuc/*.yaml}), leaving 117 cells.
Results are in \pth{experiments/results/t8/uiuc_panel_summary.json},
regenerable with \texttt{python -m cins.benchmarks.panel\_summary}.

\begin{table}[h]
\centering
\caption{UIUC panel outcome by population. Source: experiments/results/t8/uiuc\_panel\_summary.json; experiments/results/t8/ANALYSIS.md, ``H1 addendum, part 2.''}
\label{tab:uiuc-outcome}
\footnotesize
\begin{tabularx}{\textwidth}{@{}l L c@{}}
\toprule
Population & Definition & Count \\
\midrule
Cells & Sections with a generated configuration & 117 \\
No inverse problem posed & Target solve did not converge, so no target pressure distribution exists & 24 \\
Attempted & Cells where a target existed and the inverse was run & 93 \\
Converged & Newton reached the residual tolerance & 83 \\
Attempt failures & Solver or geometry faults during the attempt & 10 \\
\bottomrule
\end{tabularx}
\end{table}

\paragraph{Result: accuracy holds on every converged section; the composite
criterion fails on iteration count.} Of the 117 cells, 24 pose no inverse
problem because the target solve itself does not converge, so no target
pressure distribution exists to match. Of the remaining 93 attempted
cells, 83 converge and 10 fail during the attempt. All 83 converged cells
recover the free coefficients to $\mathrm{err\_free\_inf}<10^{-4}$, a
1.000 success fraction with Wilson 95\% lower bound 0.956. The
iteration-budget criterion, $\leq9$ Newton iterations, holds for 73 of the
83 converged cells (0.880, Wilson LB 0.792). The pre-registered composite
criterion, both accuracy and iteration budget together, is 73/83 = 0.880
among converged cells (Wilson LB 0.792) and 73/93 = 0.785 among attempted
cells (Wilson LB 0.691), both below the pre-registered 0.9 threshold.

\begin{table}[h]
\centering
\caption{UIUC panel criteria and Wilson 95\% lower bounds. Source: experiments/results/t8/uiuc\_panel\_summary.json.}
\label{tab:uiuc-criteria}
\begin{tabular}{lcc}
\toprule
Criterion & Result & Wilson 95\% lower bound \\
\midrule
Coefficient accuracy ($\mathrm{err\_free\_inf}<10^{-4}$), among converged & 83/83 = 1.000 & 0.956 \\
Iteration budget ($\leq9$ iterations), among converged & 73/83 = 0.880 & 0.792 \\
Both criteria, among converged & 73/83 = 0.880 & 0.792 \\
Both criteria, among attempted & 73/93 = 0.785 & 0.691 \\
\bottomrule
\end{tabular}
\end{table}

Ten converged sections exceed the iteration budget, requiring between 10
and 50 Newton iterations while still recovering coefficients to between
$1\times10^{-11}$ and $3\times10^{-10}$, the same accuracy band as the
sections that converge within budget. The single-digit iteration
expectation in the dossier was calibrated on NACA sections
(\S\ref{sec:h1-panel}). It does not transfer to the wider UIUC geometry
space, where sections with unusual camber or thickness distributions
converge more slowly without loss of final accuracy.

\paragraph{Exclusion classes.} Table~\ref{tab:uiuc-exclusions} lists the
exclusion classes spanning the 24 not-posed and 10 attempt-failure cells.

\begin{table}[h]
\centering
\caption{UIUC panel exclusion classes. Source: experiments/results/t8/uiuc\_panel\_summary.json.}
\label{tab:uiuc-exclusions}
\begin{tabular}{lc}
\toprule
Exclusion class & Count \\
\midrule
\texttt{target\_natural} & 16 \\
\texttt{target\_forced} & 7 \\
\texttt{solver\_stagnation} & 3 \\
\texttt{other} & 3 \\
\texttt{geometry\_te\_gap} & 3 \\
\texttt{other:RuntimeError} & 1 \\
\texttt{init\_flow\_solve} & 1 \\
\bottomrule
\end{tabular}
\end{table}

The \texttt{target\_natural} and \texttt{target\_forced} classes account
for 23 of the 24 not-posed cells: the flow solver cannot produce a
converged solution for the section itself, so no target exists and the
inverse formulation is not exercised. The \texttt{geometry\_te\_gap} class
covers three sections, \texttt{goe780}, \texttt{mh62}, and
\texttt{naca64206}, that fail mfoil's wake-direction assertion
(\pth{vendor/mfoil/mfoil.py}, \texttt{build\_wake}) after the
minimum-trailing-edge-gap fix that resolved the sharp-trailing-edge case.
Their trailing-edge gap is already above the minimum, so the cause is
distinct from the case the fix addressed. This is recorded as an open item
in \S\ref{sec:limitations} rather than resolved here.

\paragraph{Combined H1 verdict (NACA + UIUC).} Accuracy holds on every
section that converged across both panels: 18 of 18 eligible NACA sections
and 83 of 83 converged UIUC sections recover the free coefficients to
$\mathrm{err\_free\_inf}<10^{-4}$, with observed errors clustering between
$10^{-11}$ and $10^{-9}$. The pre-registered composite criterion, which
also required the solve to complete within 9 Newton iterations, is not met
on the UIUC panel: the Wilson lower bound on the composite criterion is
0.792 among converged UIUC cells and 0.691 among attempted UIUC cells,
both below the pre-registered 0.9 threshold. The cause is the iteration
budget, not accuracy. The single-digit iteration expectation in the
dossier was calibrated on NACA sections and does not transfer to the wider
UIUC geometry space, where sections with unusual camber or thickness
distributions take longer to converge while still reaching the same
accuracy. \textbf{H1 verdict:} the hypothesis that the monolithic
formulation recovers geometry on realisable targets is supported on
accuracy across both panels, 18/18 NACA sections and 83/83 converged UIUC
sections out of a 117-section UIUC panel. The pre-registered composite
criterion, which also required single-digit iterations, is not met on the
UIUC panel. The failure is isolated to the iteration-count component; the
accuracy component holds on every converged section in both panels.

\subsection{Flow-solve count vs.\ nested optimisation}
\label{sec:h2}

\paragraph{The corrected headline.} STATS\_PROTOCOL pre-registered a
$\geq$100x flow-solve-reduction claim; it is \textbf{rejected} under both
fair-paired controls actually run. Table~\ref{tab:h2} pairs each monolithic
run against a nested Levenberg--Marquardt control sharing the same
initialisation strategy, so the comparison is apples-to-apples in the same
currency (\texttt{n\_flow\_solves\_equivalent}, both sides).

\begin{table}[h]
\centering
\caption{Flow-solve-count comparison, monolithic vs.\ nested LM control. Source: paper/results\_digest.md~\S5.3; experiments/results/t8/ANALYSIS.md~\S2.}
\label{tab:h2}
\begin{tabular}{lccc}
\toprule
Pairing (shared init strategy) & Monolithic solves & Nested LM solves & Ratio \\
\midrule
Presolve-init & 120 & 373 & $\mathbf{3.1\times}$ \\
Perturbed-init (cold-start) & 43 & 347 & $\mathbf{8.1\times}$ \\
\bottomrule
\end{tabular}
\end{table}

The nested control is \texttt{scipy.optimise.least\_squares(method="lm",
ftol=xtol=gtol=1e-8, x\_scale="custom"(\allowbreak $|x_0|$ floored at
$10^{-3}$))}~\citep{virtanen2020scipy}, converged both times (\texttt{xtol}
termination), over the same 16 free coefficients. This is a competent
modern LM solver, not a weak strawman: it converges the cold-start case in
\texttt{n\_nfev=6} outer LM iterations, each costing multiple flow solves
for its own finite-difference Jacobian columns. \textbf{Recommended P1
sentence} (paper/results\_digest.md~\S5.3): \emph{``The monolithic
architecture requires 3.1--8.1x fewer flow-residual evaluations than a
fairly-paired, well-tuned nested Levenberg--Marquardt baseline
(\texttt{scipy.least\_squares}), depending on initialisation strategy,
substantially less than the order-of-magnitude reduction hypothesised a
priori, though the qualitative advantages of determinism, guaranteed local
quadratic convergence, and exact analytic constraint-row Jacobians persist
independent of this ratio.''} Figure~\ref{fig:h2}.

\begin{figure}[h]
\centering
\includegraphics[width=0.75\textwidth]{h2_flow_solves.png}
\caption{Flow-solve-count comparison, linear non-truncated axis, both fair-paired comparisons, pre-registered 100x line shown for scale. Source: \pth{experiments/results/t8/figures/h2_flow_solves.png}.}
\label{fig:h2}
\end{figure}

\subsection{Presolve vs.\ station selection as the uniqueness guard}

\texttt{t8\_init\_perturbed} (no presolve, \texttt{station\_selection=
qr\_pivot}): 43 flow-solve-equivalents, 5 iterations, error
$2.12\times10^{-11}$, submap cond 89.8, essentially identical to the
baseline's 90.3, since station selection is a property of the stations
(computed once from the T4 sensitivity matrix), not of the initial guess.
\texttt{t8\_n08\_baseline} (presolve, same station selection): 120
flow-solve-equivalents, 4 iterations, error $1.08\times10^{-11}$, submap
cond 90.3. Interpretation: skipping the presolve costs one extra Newton
iteration but \emph{reduces} total flow-solve count ($43<120$) and does not
cost correctness. The presolve consumes roughly 77 flow-solve-
equivalents ($\approx$64\% of the baseline total), primarily for iteration
efficiency and starting-residual quality, not for the uniqueness guarantee,
which comes from \texttt{station\_selection=qr\_pivot} alone
(paper/results\_digest.md~\S5.4). This has \emph{not} been verified in the
sharper crossed form (\texttt{init=perturbed} $\times$
\texttt{station\_selection=even}); it is flagged explicitly as future work,
not as tested (experiments/results/t8/ANALYSIS.md~\S5, ``Caveat'').

\subsection{Quadratic convergence tail}
\label{sec:h3}

STATS\_PROTOCOL's 3-point local order estimator
($p\approx\log(R_k/R_{k-1})/\log(R_{k-1}/R_{k-2})$) is applied only where
$\geq3$ residual values exist strictly above the $\approx10^{-9}$
FD-noise floor. The raw \texttt{convergence\_order} field stored per cell
does \emph{not} apply this exclusion and is not quoted directly
(paper/results\_digest.md~\S5.5). Table~\ref{tab:h3} reports the recomputed
orders for the five cells with enough clean-tail residuals to estimate one.

\begin{table}[h]
\centering
\caption{Recomputed convergence order, clean (floor-excluded) tail. Source: paper/results\_digest.md~\S5.5; experiments/results/t8/ANALYSIS.md~\S3.}
\label{tab:h3}
\begin{tabular}{lc}
\toprule
Cell & Recomputed order (clean tail) \\
\midrule
\texttt{t8\_alpha\_free} & 1.09 \\
\texttt{t8\_init\_perturbed} & 2.08 \\
\texttt{t8\_n06} & 2.46 \\
\texttt{t8\_transition\_free} & 2.15 \\
\texttt{t8\_n12} & 4.24 (late-tail 3-point estimate; read as ``clearly superlinear,'' not literal quartic order) \\
\bottomrule
\end{tabular}
\end{table}

\textbf{Median order: 2.15} ($n{=}5$ estimable cells), meeting the pre-
registered $\geq1.8$ target. The remaining 7/12 converged cells cannot be
scored by this method because they reach the solver's residual floor in
only 2 post-initial steps, too fast for a 3-point estimator to have 3 clean
points. This is itself consistent with (though it does not formally
confirm) fast superlinear-to-quadratic convergence: the baseline, for
example, drops from $O(10^{-4})$ to the $\sim10^{-10}$ floor in exactly 2
Newton steps, a $\sim$150x-per-step reduction rate not consistent with
linear convergence (experiments/results/t8/ANALYSIS.md~\S3).

\subsection{Comparative positioning against other inverse-design approaches}
\label{sec:comparison}

Tables~\ref{tab:comparison-quant} and~\ref{tab:comparison-struct} place this
paper's architecture against five other approaches to the same underlying
problem (given a target surface pressure or speed distribution, recover a
geometry that produces it), each representing a structurally distinct
formulation class. The comparison is split across two tables because its two
halves are not commensurable in the same way:
Table~\ref{tab:comparison-quant} collects the quantitative axes, where the
cost column is reported in whatever unit each formulation class naturally
exposes, and Table~\ref{tab:comparison-struct} collects the structural
guarantees, where the classes separate sharply and no numbers are involved at
all. Rows (i) and (ii) are this paper's own measured, fair-paired numbers
(\S\ref{sec:h2}); rows (iii) to (vi) are drawn from the cited literature and
are explicitly marked \emph{qualitative} wherever the source does not report a
directly comparable metric. No number in either table is asserted unless it is
a repository artifact from rows (i) or (ii), or attributed to a specific cited
source.

\begin{table}[p]
\centering
\caption{Comparative positioning, quantitative axes: formulation class, cost
per design, and the precision each approach has been demonstrated to. Cost is
reported in the unit each formulation class naturally exposes, which is why
the column is not a single commensurable number; ``qual.'' marks an entry for
which no directly comparable published metric exists.}
\label{tab:comparison-quant}
\footnotesize
\begin{tabularx}{\textwidth}{@{}>{\raggedright\arraybackslash}p{2.45cm} L L L@{}}
\toprule
Method & Formulation class & Flow solves per design & Demonstrated precision \\
\midrule
(i) \textbf{This work} (CINS monolithic, \S\ref{sec:formulation}) &
  constrained root-finding, single square Newton system &
  \textbf{43--120} measured, T7 winning configuration, two initialisation
  strategies (\S\ref{sec:h2}) &
  $\norm{A-A^*}_\infty$ at the $10^{-11}$ level on self-consistency
  (Table~\ref{tab:t7}) \\
(ii) Nested LM baseline (\pth{scipy.least_squares}, \S\ref{sec:h2}) &
  outer-loop nonlinear least-squares optimisation &
  \textbf{347--373} measured on the same paired cases (\S\ref{sec:h2}) &
  converges to \texttt{xtol} against the same coefficient-level target,
  paired with row (i) \\
(iii) MISES Modal-Inverse~\citep{drela1990viscous} &
  constrained root-finding, mode-shape amplitudes appended to the same global
  Newton system &
  one Newton solve per design, architecturally identical to row (i); exact
  count not published &
  qual.: production-validated on transonic cascades over decades, no single
  benchmark number available for direct comparison \\
(iv) CST-smoothing inverse~\citep{lane2010inverse} &
  iterative outer loop; CST used as a \emph{post-hoc smoother} on a
  pressure-residual-driven shape correction, solver-agnostic &
  qual.: multiple outer iterations, count depends on the wrapped solver and is
  not reported as a fixed number in the source &
  qual.: subsonic and transonic cases shown converging; no
  coefficient-recovery benchmark reported \\
(v) Learned generative inverse (cGAN~\citep{yilmaz2020cgan}; CDDPM
\citep{yu2025cddpm}; LF-PINN~\citep{lfpinn2024}) &
  amortised statistical inference, offline-trained generative or surrogate map &
  approximately zero \emph{online} per design after training; training-set
  generation requires $O(10^3\text{--}10^5)$ offline flow solves, typical for
  this literature &
  CDDPM reports a 34.4 percent precision improvement over WGAN baselines
 ~\citep{yu2025cddpm}, a relative rather than absolute geometric-recovery
  metric; LF-PINN reports at least a tenfold online speedup against direct
  mfoil analysis~\citep{lfpinn2024} \\
(vi) Adjoint target-$C_p$ optimisation (SU2~\citep{economon2016su2}) &
  gradient-based PDE-constrained optimisation, primal and adjoint pairs per
  design-cycle iteration &
  qual.: $O(10\text{--}100)$ primal/adjoint pairs typical for shape
  optimisation of this scale, case- and tolerance-dependent &
  qual.: mature and widely validated on shape-optimisation benchmarks; not run
  in this study, so no head-to-head number is available \\
\bottomrule
\end{tabularx}
\end{table}

\begin{table}[p]
\centering
\caption{Comparative positioning, structural axes: what each approach
guarantees about the solution it returns, how it handles geometric design
constraints, and its availability. The determinism column is the one on which
the formulation classes separate most sharply.}
\label{tab:comparison-struct}
\footnotesize
\begin{tabularx}{\textwidth}{@{}>{\raggedright\arraybackslash}p{2.45cm} L L L@{}}
\toprule
Method & Determinism and uniqueness guard & Constraint handling &
Licence and availability \\
\midrule
(i) \textbf{This work} (CINS monolithic) &
  deterministic Newton; QR-pivoted station selection is an explicit, measured
  identifiability guard (\S\ref{sec:ablations}) &
  exact linear equality rows $G A = b$, with a zero-cost constant Jacobian
  block &
  open~\citep{cinsrepo}; mfoil is MIT-licenced
 ~\citep{fidkowski2022coupled} \\
(ii) Nested LM baseline &
  deterministic given fixed initialisation and tolerances; no analogous
  uniqueness guard, relying on local LM convergence only &
  none built in; bounds and scaling only, with no analytic constraint rows in
  this configuration &
  open~\citep{virtanen2020scipy} \\
(iii) MISES Modal-Inverse &
  deterministic Newton; no published uniqueness or identifiability study of
  mode count against target sampling analogous to \S\ref{sec:ablations} &
  mode amplitudes are heuristic bump functions rather than analytic geometric
  functionals (\S\ref{sec:related}) &
  gated: free for faculty-signed academic use, commercial licence otherwise
 ~\citep{drela1990viscous}; no open-source port exists \\
(iv) CST-smoothing inverse &
  deterministic given fixed settings; no formal uniqueness guard reported &
  no analytic constraint rows; the correction is smoothed but not constrained
  to a design-parameter manifold &
  method description only, solver-agnostic, with no released code cited \\
(v) Learned generative inverse &
  \emph{not} deterministic for GAN and diffusion variants, which sample
  stochastically; LF-PINN's surrogate is deterministic but is an
  approximation rather than an exact solve. \textbf{No convergence guarantee
  or exact-recovery certificate exists for any variant.} &
  none intrinsic to generation; plausible-geometry constraints are learned
  from the training distribution rather than enforced algebraically &
  open-source code released with most cited papers, but a training database is
  required \\
(vi) Adjoint target-$C_p$ optimisation &
  deterministic gradient descent to a \emph{local} optimum; no uniqueness
  guard, which is the sense in which adjoint target-$C_p$ matching remains
  optimisation rather than root-finding \dossierref{\S6.3} &
  general nonlinear constraints via adjoint sensitivities: flexible, but
  iterative rather than closed-form &
  open source, LGPL~\citep{economon2016su2} \\
\bottomrule
\end{tabularx}
\end{table}

Three patterns follow from the table. First, only rows (i) and (iii) are
architecturally root-finding rather than optimisation or amortised
inference. This distinction, argued in this paper's Introduction
(\S\ref{sec:related}), is not cosmetic: a converged root-find has a
residual-norm certificate that a converged optimisation or a sampled
generative output does not. Second, the flow-solve-count column spans
roughly four orders of magnitude across the table, not because the methods
disagree about a shared quantity, but because ``flow solve'' means a
different thing in each formulation class (a Newton linearisation vs.\ an
LM outer iteration vs.\ an offline training epoch vs.\ a primal/adjoint
pair). This is exactly the currency-matching discipline \S\ref{sec:h2}
enforces for the one comparison (row (i) vs.\ row (ii)) this paper actually
measured; rows (iii)--(vi) are reported qualitatively for this reason, not
omitted. Third, the licence/availability column is not incidental context:
it is the practical reason MISES's built-in Modal-Inverse mode (row (iii)),
despite being architecturally identical to this paper's contribution, has
not seen the open comparative study this table attempts. A \$10k-gated,
non-open-source solver is not reproducible science. This is also why mfoil
\citep{fidkowski2022coupled} rather than MISES is this paper's Stage 1
testbed (\S\ref{sec:formulation}), and why the Stage 3 extension
(\S\ref{sec:conclusions}) targets porting the architecture \emph{into}
MISES rather than requiring readers to already have a MISES licence to
reproduce this paper's own results.
